Una partícula amb una càrrega $q = -1,60\times 10^{-19}\ \mathrm{C}$ i una massa $m = 1,70\times 10^{-27}\ \mathrm{kg}$ entra amb una velocitat $\vec{v} = v\,\vec{\imath}$ en una regió de l'espai en la qual hi ha un camp magnètic uniforme $\vec{B} = -0,50\ \mathrm{T}\ \vec{k}$. El radi de la trajectòria circular que descriu és $r = 0,30\ \mathrm{m}$.

\figura[0.35]{fig1.png}

\begin{apartat}{1}
Dibuixeu la força que fa el camp sobre la partícula en l'instant inicial i calculeu la velocitat $v$.
\end{apartat}

\begin{apartat}{1}
Calculeu el període del moviment i la velocitat angular. Calculeu l'energia cinètica de la partícula en el moment que entra en el camp magnètic i també després de fer una volta completa.
\end{apartat}
